\section{张量价格与连线距离}\label{sec:prices-background}

一个秩为 \(r\) 的双线性算法通过局部线性映射，从 \(r\) 个独立的标量乘积构造出其张量。\emph{谱点}给张量赋予一个值，在这类映射下，该值只会减小或保持不变。取张量积使这些值相乘，从而使它们的对数相加。我们称这些对数为\emph{价格}。对于连线，价格衡量在其两个站点之间共享一个指标的成本。谱点正是以这种方式赋予张量网络的站点以几何结构。

\subsection{从张量到距离}

固定一个域和一个至少含两个站点的有限集合 \(S\)。每个张量在每个站点上对应一个有限维向量空间。如果可以在各个站点分别选取一个线性映射，使它们共同把 \(U\) 映到 \(T\)，就称 \(T\) 为 \(U\) 的一个\emph{限制}。取张量积时，将各个站点上的因子分别归组。取直和时，在每个站点上取向量空间的直和，使两个张量处于分开的块中。单位张量 \(\mathbf 1\) 是各空间均为一维时的张量 \(1\)。

谱点 \(\varphi\) 是一个取值于 \(\mathbb R_{\ge0}\)、在限制下单调的映射，它保持零和单位，且满足
\[
     \varphi(T\oplus U)=\varphi(T)+\varphi(U),\qquad
     \varphi(T\otimes U)=\varphi(T)\varphi(U).
\]
每个非零张量都可限制为 \(\mathbf 1\)，故其取值至少为一。在尺度 \(D>0\) 下，价格为 \(H(T)=D\log\varphi(T)\)，并约定 \(H(0)=-\infty\)。本节与第一部分使用自然对数。因此在相同尺度下，本文的价格是文献~\cite{oaG} 中价格的 \(\log 2\) 倍。

对非空的 \(A\subseteq S\)，张量
\[
         \Delta_A(q)=\sum_{j=1}^q\bigotimes_{i\in A} e_j
\]
在 \(A\) 中各站点之间共享一个可取 \(q\) 个值的指标。它的其他腿均为一维。两个对角张量的张量积使其维数相乘，局部投影则可以丢弃任意选定的一组标签。因此，比较维数 \(q\) 与维数 \(2\) 的对角张量的各次幂，就有
\[
            H(\Delta_A(q))=\delta(A)\log q,\qquad q\ge2.
\]
我们称 \(\delta(A)\) 为 \(A\) 上的共享标签费率。特别地，\(\Delta_S(q)\) 是 \(q\) 个单位张量的直和，故 \(\delta(S)=D\)。对于 \(a\ne b\)，\(\Delta_{\{a,b\}}(q)\) 是一条维数为 \(q\) 的连线。记 \(d(a,b)=\delta(\{a,b\})\)，并令 \(d(a,a)=0\)。例如，按 \(S\) 的一个固定切分展平所得的矩阵秩是一个谱点。在尺度一处，分居切分两侧的站点之间的连线距离为一，同侧站点之间的连线距离为零。

在两条连线的公共端点处作收缩，便得到连接另外两个端点的一条连线。价格的单调性与可加性给出 \(d\) 的三角不等式。一条 \(q_i\) 维连线还可将一条维数为 \(q_i\) 的腿从 \(x_i\) 移动至 \(y_i\)。移动非零张量 \(T\) 的每一条腿，再沿反方向应用相同论证，便得运输界
\begin{equation}\label{eq:prices-transport}
       |H(T;x)-H(T;y)|\le\sum_i d(x_i,y_i)\log q_i .
\end{equation}
其中 \(x=(x_i)\) 和 \(y=(y_i)\) 表示各条腿的放置位置。放在同一站点的腿归为一组。当所有腿都位于一个站点时，该处的任一非零张量都与单位张量等价，价格为零。

\subsection{几何扩张与谱比较}

\emph{谱空间}是一个度量空间，在其一切有限放置上都有相容的价格，且连线费率就是给定的距离。相容性允许添入平凡的腿，也允许将位于重合站点的腿归组。在直径为正的有限集合上，尺度等于其共享标签费率，可以随集合而变化。在单个点处，尺度为零。把连线距离为零的站点等同，就可将有限多个站点上的谱价格化为一个谱空间。

每个谱空间连同它的全部价格都可嵌入一个完备的 \(\operatorname{CAT}(0)\) 谱空间。在这一扩张中，每对点 \(a,b\) 有唯一中点 \(m\)，且每个点 \(o\) 都满足
\[
      d(o,m)^2\le \tfrac12 d(o,a)^2+\tfrac12 d(o,b)^2
                                      -\tfrac14 d(a,b)^2 .
\]
当各条腿沿测地线移动时，价格是联合凸的。距离与价格可以按同一比例缩放。在这类空间的带基点度量超极限中，一个固定张量在有界放置序列上的价格，等于其各个价格的超极限~\citep[第~\ext{G}{sec:spectra}、\ext{G}{sec:limits} 节，定理~\ext{G}{thm:geometry} 与命题~\ext{G}{prop:joint-convexity}]{oaG}。我们还用到距离的凸性，以及到非空闭凸集的最近点投影的非扩张性。一个具有有限支撑的概率测度存在唯一的重心 \(b\)，使其平均平方距离最小。在 \(z\) 处计算同一平均值，其增量至少为 \(d(z,b)^2\)~\citep{BridsonHaefliger,oaG}。

在代数闭域上，固定张量各因子的维数后，价格的次水平集是 Zariski 闭的。作域扩张时，谱点在一组定义于原域的张量上所能同时取得的值保持不变~\citep[第~\ext{G}{sec:spectra} 节]{oaG}。最后，设同一组站点上的两个非零张量对每个谱点均满足 \(\varphi(T)<\varphi(U)\)，则存在一个固定整数 \(C\)，使得对于每个 \(j\ge1\)，\(T^{\otimes j}\) 都是 \(C\) 份 \(U^{\otimes j}\) 的直和的一个限制。我们用这种严格的谱比较形式把价格不等式转化为精确的张量限制~\citep[引理~\ext{G}{lem:spectral-comparison}]{oaG}。
